\section{System overview}
\label{sec:overview}

This section sketches the system in just enough detail for the math that
follows. It assumes familiarity with Royco-style tranched markets at the
level of Dawn, and with Balancer V3's pool, hook, and rate-provider
architecture.

\subsection{Tranches and the loss waterfall}
\label{sec:overview:tranches}

A Royco market splits deposits across two tranches with different risk
profiles:

\begin{description}
  \item[Senior (ST)] Protected capital. Earns a fixed share of yield;
    absorbs losses only after JT's buffer is exhausted.
  \item[Junior (JT)] First-loss coverage capital. Absorbs senior losses up
    to its effective NAV; in exchange earns yield share and coverage
    premiums.
\end{description}

Each tranche has two NAV figures the accountant tracks per sync:

\begin{description}
  \item[Raw NAV] The mark-to-market value of the tranche's \emph{own}
    underlying assets (\code{ST_RAW_NAV}, \code{JT_RAW_NAV}).
  \item[Effective NAV] The value each tranche \emph{controls} after the
    waterfall distributes PnL according to coverage rules (\code{ST_EFF},
    \code{JT_EFF}).
\end{description}

The waterfall, encoded in
\code{RoycoAccountant.\_previewSyncTrancheAccounting}, is a fixed-priority
allocator. Given the two raw-NAV deltas since the last checkpoint, it routes
them through four branches (junior loses, junior gains, senior loses, senior
gains), each a chain of \code{min} / \code{max} / addition / subtraction.
We refer to this allocator as the function $F$ throughout this document.

Two structural facts about $F$ matter later:

\begin{enumerate}
  \item \textbf{NAV conservation:}
    $\STraw + \JTraw \equiv \STeff + \JTeff$. The waterfall only
    redistributes value between tranches; it never creates or destroys it.
    The on-chain \code{NAV_CONSERVATION_VIOLATION} guard enforces this on
    every sync.
  \item \textbf{Piecewise-affine:} each branch is affine in its inputs,
    with slope $\in \{0, 1\}$. This is the property invoked in the
    existence proof (\cref{sec:existence}).
\end{enumerate}

\subsection{What makes Dusk different from Dawn}
\label{sec:overview:dawn-vs-dusk}

In a Dawn market, the JT side deposits a plain ERC20 (e.g.\ \code{sNUSD},
\code{USDC}). The kernel custodies the token directly and reads
\code{JT_RAW_NAV} from an exogenous oracle.

In a Dusk market, the JT asset is a Balancer V3 Pool Token (BPT) --- the
LP token of a pool pairing two assets:

\begin{itemize}
  \item the senior tranche share token;
  \item a stable quote asset (\code{USDC}, \code{sUSDs}, \code{srRoyUSDC},
    etc.).
\end{itemize}

The pool plays a dual role: it serves as JT's first-loss coverage collateral
(the Dawn function) \emph{and} as a secondary-market exit venue for senior
holders who want to swap out of their tranche shares into stables (the new
Dusk function).

The consequence is that JT's collateral value now depends on the price of
the senior share, which is the rate $y = \STeff / N$, where $N$ is the
total senior share supply. But $\STeff$ is the output of the waterfall $F$,
and computing $F$ requires marking JT's BPT to market --- which requires
$y$. This is the circularity the rest of this document addresses, starting
in \cref{sec:setup}.

\subsection{Pieces of a Dusk market}
\label{sec:overview:pieces}

A Dusk market is composed of:

\begin{description}
  \item[Gyro E-CLP pool] A Balancer V3 elliptic concentrated liquidity
    pool over $\{\text{senior share},\ \text{quote asset}\}$. The senior
    share token is registered as a \code{WITH_RATE} token, so Balancer
    scales its live balance by the rate provider's published rate before
    pool math runs.
  \item[Dusk kernel] One contract wearing two hats:
    \begin{itemize}
      \item the senior token's rate provider, exposing \code{getRate()}
        which returns the cached senior share rate;
      \item the sync entrypoint, which runs the bisection solve
        (\cref{sec:solver}) and writes the freshly-solved rate to the
        cache.
    \end{itemize}
  \item[Hook contract] Forwards Balancer V3's before-operation callbacks
    (swap, add liquidity, remove liquidity) into the kernel's sync
    entrypoint, so each operation triggers a fresh solve before the pool
    math runs.
  \item[Accountant] Reused from Dawn; called by the kernel during the
    bisection solve to evaluate $F$.
\end{description}
